% Generated draft: requires mathematical and source review.
\documentclass[11pt]{article}
\usepackage[margin=25mm]{geometry}
\usepackage{fontspec}
\setmainfont{Latin Modern Roman}
\usepackage{amsmath,amssymb,mathtools,unicode-math}
\setmathfont{Latin Modern Math}
\usepackage{enumitem,needspace,xurl,hyperref}
\setlength{\parindent}{0pt}
\setlength{\parskip}{5pt}
\setlength{\emergencystretch}{3em}
\newcommand{\N}{\mathbb{N}}
\newcommand{\Z}{\mathbb{Z}}
\newcommand{\Q}{\mathbb{Q}}
\newcommand{\R}{\mathbb{R}}
\newcommand{\C}{\mathbb{C}}
\newcommand{\F}{\mathbb{F}}
\newcommand{\A}{\mathbb{A}}
\newcommand{\PP}{\mathbb{P}}
\newcommand{\E}{\mathbb{E}}
\newcommand{\eps}{\varepsilon}
\newcommand{\dd}{\,\mathrm{d}}
\begin{document}
\section*{UnsolvedMath ID: 800763}
\textbf{Source problem code:} catalogue entry
\subsection{Definitions and mathematical statement}
Let \(K\) be a field, with no restriction on its characteristic. Its multiplicative group is \(K^\times=K\setminus\{0\}\). The field \(K\) is \emph{radically closed} if \(K^\times\) is divisible, meaning that
\[
\forall n\in\Z_{\geq 1}\;\forall a\in K^\times\;\exists b\in K^\times:\quad b^n=a.
\]
A \(K\)-curve means a separated scheme of finite type over \(K\) of dimension one. A field \(K\) is \emph{ample} if every smooth geometrically integral \(K\)-curve \(C\) with a \(K\)-rational point has infinitely many \(K\)-rational points. Here geometrically integral means that \(C\times_K\overline K\) is integral, where \(\overline K\) is an algebraic closure of \(K\).

\textbf{Question (Bary-Soroker and Fehm, Question 6, 2011).} Is every radically closed field ample? Explicitly, does every field \(K\) satisfy the implication
\[
\left(\forall n\geq 1\;\forall a\in K^\times\;\exists b\in K^\times:\ b^n=a\right)
\Longrightarrow
\left(\forall C\text{ smooth geometrically integral }K\text{-curve},\quad C(K)\neq\varnothing\Longrightarrow \#C(K)=\infty\right)?
\]

\textbf{Unresolved source qualification.} The supplied background also asserts that every radically closed field contains all roots of unity and hence, by Kummer theory, has no solvable finite Galois extensions, understood as nontrivial extensions. These are retained as supplied background claims, without independent verification. Their unrestricted characteristic scope requires clarification: in positive characteristic, Kummer theory alone does not address extensions whose degree is divisible by the characteristic. These background claims are not added as hypotheses of the question.
\subsection{Short English statement}
The question asks whether, over every field in which each nonzero element has an \(n\)-th root for every positive integer \(n\), each smooth geometrically integral curve with one rational point must have infinitely many rational points. The supplied registry labels the question open; that status is not independently verified. The background claims concerning roots of unity and solvable extensions require clarification in positive characteristic.
\subsection{Sources}
\begin{itemize}
\item Supplied local UnsolvedMath registry record: permanent ID \(800763\), source problem code \texttt{catalogue entry}, title \emph{Open Problems in the Theory of Ample Fields --- Question 6}. The record attributes the question to Bary-Soroker and Fehm, gives the proposed year as 2011, and labels its status \texttt{open}.
\item Bary-Soroker and Fehm, \emph{Open Problems in the Theory of Ample Fields} (2011), Question 6. Supplied source link: \url{https://arxiv.org/abs/1106.1310}.
\item Supplied provenance: accessed 2026-10-16; extraction method \texttt{arxiv\_latex}. This access date and provenance are supplied record metadata, not a claim of independent source access or verification by this draft.
\end{itemize}
\end{document}
